%------------------------------------------------------------------------------%
\section{Integrais Impróprias}
\label{sec:Integrais_Improprias}

As integrais definidas são calculadas em um intervalo finito e a função $f$
não deve ter nenhuma descontinuidade infinita. Porém, mesmo se o intervalo
for infinito ou $f$ possuir uma descontinuidade infinita, em alguns casos,
é possível calcular a área sob o gráfico usando a integral imprópria.

%------------------------------------------------------------------------------%
\subsection{Integrais impróprias do tipo I -- Intervalo de integração ilimitado}

%------------------------------------------------------------------------------%
\begin{definicao}{Integral Imprópria do tipo I}{def:intergal-impropria}
  Seja $f$ função contínua em intervalo ilimitado, do tipo
  $[a, \infty)$,
  $(-\infty, b]$ ou
  $(-\infty, \infty)$.
  A integral definida de $f$ em cada
  um desses intervalos é dita imprópria e é dada respectivamente por
  \begin{align*}
    \int_a^{\infty} f(x)                  \, dx
    & = \lim_{t\to\infty} \int_a^t f(x)   \, dx \\[3mm]
        \int_{-\infty}^b f(x)             \, dx
    & = \lim_{t\to-\infty} \int_t^b f(x)  \, dx \\[3mm]
        \int_{-\infty}^{\infty} f(x)      \, dx
    & = \lim_{t\to-\infty} \int_t^c f(x)  \, dx
      + \lim_{t \to \infty} \int_c^t f(x) \, dx,
        \quad \text{para algum } c \in \R.
  \end{align*}
  Dizemos que a integral imprópria converge se o(s) limite(s) existe(m);
  caso contrário, diverge.
\end{definicao}
%------------------------------------------------------------------------------%

Se \(\int_a^{\infty} f(x) dx\)
converge, então \(\lim_{x \to \infty} f(x) = 0\).

%------------------------------------------------------------------------------%
\begin{example}
  Determine se a integral \(\int_1^{\infty} \frac{1}{x} dx\)
  é convergente ou divergente.
  \solution
  De acordo com a definição de integrais impróprias do tipo $I$, temos
  \begin{align*}
   \int_1^{\infty}  \frac{1}{x} dx
    & = \lim_{t \to \infty} \int_1^{t} \frac{1}{x} dx   \\[3mm]
    & = \lim_{t \to \infty} \ln\abs{x} \evalat{1}{t}    \\[3mm]
    & = \lim_{t \to \infty} \big(\ln(t)- \ln(1)\big)    \\[3mm]
    & = \lim_{t \to \infty} \ln(t)= \infty.
  \end{align*}
  Portanto, a integral imprópria \(\int_1^{\infty} \frac{1}{x}dx\)
  é divergente.
\end{example}
%------------------------------------------------------------------------------%

%------------------------------------------------------------------------------%
\begin{example}
  Determine se a integral \(\int_{-\infty}^0 2^x dx\)
  é convergente ou divergente.
  \solution
  Pela definição de integral imprópria
  \begin{align*}
    \int_{-\infty}^0 2^x dx
    & = \lim_{t \to - \infty} \int_t^{0} 2^x dx \\
    & = \lim_{t \to - \infty} \frac{2^x}{\ln(2)} \evalat{t}{0}  \\
    & = \lim_{t \to - \infty} \frac{1}{\ln(2)} \left( 2^0 - 2^t \right) \\
    & = \frac{1}{\ln(2)}
  \end{align*}
  Logo a integral imprópria \(\int_{-\infty}^0 2^x dx\) é convergente e vale
  \(\frac{1}{\ln(2)}\).
\end{example}
%------------------------------------------------------------------------------%

%------------------------------------------------------------------------------%
%\begin{example}
%  Determine se a integral \(\int_{-\infty}^{\infty} x^3e^{-x^4} dx\)
%  é convergente ou divergente.
%  \solution
%  Usando  definição de integral imprópria
% \begin{equation}
%  \label{exemplo_eqimp}
%   \int_{-\infty}^{\infty} x^3e^{-x^4} dx
%= \lim_{t \to-\infty} \int_{t}^{0} x^3e^{-x^4} dx
%   + \lim_{t \to \infty} \int_{0}^{t} x^3e^{-x^4} dx
%\end{equation}
%Para resolvermos a integral indefinida $\int x^3e^{-x^4} dx$,
%usamos a substituição $u=-x^4 \Rightarrow du=-4x^3dx$. Logo,
%\[
% \int x^3e^{-x^4} dx
% = -\frac{1}{4}\int e^u du
% = -\frac{1}{4}e^u+C
% = -\frac{1}{4}e^{-x^4}+C
%\]
% \begin{comment}
%  Usando o Teorema Fundamental do cálculo temos
%  \begin{align*}
%    \int_{-\infty}^{0} x^3e^{-x^4} dx
%    & = \lim_{t \to -\infty} \int_{t}^{0} x^3e^{-x^4} dx \\
%    & = \lim_{t \to -\infty} -\frac{1}{4}\left( e^0-e^t \right) \\
%    & =-\frac{1}{4}
%  \end{align*}
%  \begin{align*}
%    \int_{0}^{\infty} x^3e^{-x^4} dx
%    & = \lim_{t \to \infty} \int_{0}^{t} x^3e^{-x^4} dx \\
%    & = \lim_{t \to -\infty} -\frac{1}{4}\left( e^t-e^0 \right) \\
%    & =-\infty
%  \end{align*}
%  Como uma das integrais do lado direito da equação
%  \ref{exemplo_eqimp} diverge então
%  \(\int_{-\infty}^{\infty} x^3e^{-x^4} dx\)
%  é divergente.
%%\end{example}
%\end{comment}
%------------------------------------------------------------------------------%

%------------------------------------------------------------------------------%
\begin{example}
  Determine se a integral
  \(\int_{-\infty}^{\infty} \frac{1}{1+x^2} dx\)
  é convergente ou divergente.
  \solution
  Mais uma vez iremos usar a definição de integral imprópria do tipo $I$.
  Nessa definição é conveniente escolher $a=0$, assim
  \[
    \int_{-\infty}^{\infty} \frac{1}{1+x^2} dx
    = \int_{-\infty}^{0} \frac{1}{1+x^2} dx
    + \int_{0}^{\infty}  \frac{1}{1+x^2} dx
  \]
  É necessário calcular as integrais do lado direito separadamente
  \begin{align*}
    \int_{0}^{\infty} \frac{1}{1+x^2} dx
    & = \lim_{t \to \infty} \int_0^t \frac{dx}{1+x^2} \\
    & = \lim_{t \to \infty} \arctan(x) \evalat{0}{t} \\
    & = \lim_{t \to \infty} \big(\arctan(t)- \arctan(0)\big) \\
    & = \lim_{t \to \infty} \arctan(t) \\
    & = \frac{\pi}{2}.
  \end{align*}
  \begin{align*}
    \int_{-\infty}^{0} \frac{1-}{1+x^2} dx
    & = \lim_{t \to  \infty} \int_t^0 \frac{dx}{1+x^2} \\
    & = \lim_{t \to  \infty} \arctan(x) \evalat{t}{0} \\
    & = \lim_{t \to -\infty} \big(\arctan(0)- \arctan(t)\big) \\
    & = \lim_{t \to  \infty} 0-\left(-\frac{\pi}{2}\right) \\
    & = \frac{\pi}{2}.
  \end{align*}
  Como ambas as integrais são convergentes, a integral dada é convergente
  \[
    \int_{-\infty}^{\infty} \frac{1}{1+x^2} dx
    = \frac{\pi}{2} + \frac{\pi}{2}
    = \pi
  \]
\end{example}
%------------------------------------------------------------------------------%

Como \(\frac{1}{1+x^2}>0\), a integral imprópria pode ser interpretada
como a área da região infinita sob a curva \(y=\frac{1}{1+x^2}\)
e acima do eixo $x$, conforme mostra a Figura~\ref{fig:integral_impropria}.

\begin{figure}
  \centering
  \includegraphics[width=0.5\linewidth]{integral_impropria}
  \caption{Integral Imprópria.}
  \label{fig:integral_impropria}
\end{figure}

A próxima definição é outro tipo de integral imprópria,
que envolve descontinuidade.

%------------------------------------------------------------------------------%
\subsection{Integrais impróprias do tipo II -- Descontinuidade em um extremo}

%------------------------------------------------------------------------------%
\begin{definicao}{Integral Imprópria -- Tipo 2}{def:integral-impropria-2}
\begin{enumerate}

\item Se $f$ é contínua em $[a, b)$ e descontínua em $b$,
  temos que
  \[
  \int_a^b f(x)\,dx = \lim_{t\to b^{-}} \int_a^t f(x)\,dx
  \]
  Se esse limite existir (como número finito).

\item   Se $f$ é contínua em $(a, b]$ e descontínua em $a$,
  temos que
  \[
    \int_a^b f(x)\,dx = \lim_{t\to a^{+}} \int_t^b f(x)\,dx
  \]
  Se esse limite existir (como número finito).

  Nesse caso, dizemos que a integral imprópria $\int_a^b f(x)\,dx$ é
  \textbf{convergente} se o limite correspondente existir e \textbf{divergente}
  caso contrário.
\item  Se $f$ é uma função definida em $[a, b]$, contínua exceto em um ponto
  $c \in [a, b]$ e $c$ é uma descontinuidade de $f$, então definimos
  \[
    \int_a^b f(x) \, dx = \int_a^c f(x) \, dx + \int_c^b f(x) \, dx
  \]
  caso \(\int_a^c f(x)\,dx\) e \(\int_c^b f(x)\,dx\) forem convergentes.
  \end{enumerate}
\end{definicao}

%----------------------------------------------------------------------------------
Antes de fazermos os exemplos, vamos fazer algumas observações,
que nos auxiliarão nos exemplos que seguem:
\begin{enumerate}
  \item Seja $F$ a primitiva mais geral de $f$, ou seja,
  \[
    \int f(x)dx= F(x)+c,
  \]
    então, se $k$ e $b$ são constantes, vale que
  \[
    \int f(kx) dx= \frac{F(kx)}{k}+C.
  \]
  \item Além disso, \(\int f(kx+b)dx= \frac{F(kx+b)}{k}+C\). \\
  Para ficar mais claro, observe os exemplos a seguir:

  \item \(\int \cos(5x) dx= \frac{\sin(5x)}{5}+C\).

  \item \(\int e^{8x}= \frac{e^{8x}}{8}+C\).

  \item \(\int \frac{1}{x+3}dx= \ln\abs{x+3}+C\).

\end{enumerate}
Essas observações podem ser demonstradas usando substituição simples,
técnica de integração que veremos posteriormente.

\begin{example}
  Encontre, caso exista
  \(\int_{2}^5 \frac{1}{\sqrt{x-2}} dx\).
  \solution
  Note que, a integral
  \(\int_{2}^5 \frac{1}{\sqrt{x-2}} dx\) é imprópria tipo-2, pois
  \(f(x) =\frac{1}{\sqrt{x-2}} dx\) tem uma assíntota vertical em $x=2$.
  Como essa descontinuidade ocorre no extremo esquerdo do intervalo $[2, 5]$,
  usaremos a parte 2 da Definição~\ref{def:integral-impropria-2}. Desse modo
  \begin{align*}
    \int_{2}^5 \frac{1}{\sqrt{x-2}} dx
    & = \lim_{t \to 2^{+}} \int_{t}^5 \frac{1}{\sqrt{x-2}} dx  \\
    &\stackTFCD \lim_{t \to 2^{+}} \left[ 2 \sqrt{x-2} \right]_t^5 \\
    & = \lim_{t \to 2^{+}}2(\sqrt{3}-\sqrt{t-2})= 2\sqrt{3}
  \end{align*}
  Portanto, a integral imprópria é convergente. Uma vez que o integrando
  é uma função positiva, podemos interpretar o valor da integral como a
  área da região abaixo da curva $f(x)=\frac{1}{\sqrt{x-2}}$,
  no intervalo $(2, 5]$.

  Sabendo que
  \(\int \frac{1}{\sqrt{x}}dx= \int x^{-\frac{1}{2}}dx = 2\sqrt{x}+k\),
  pela observação feita anteriormente, usamos que
  \(\int \frac{1}{\sqrt{x-2}}dx= 2\sqrt{x-2}+k\).
\end{example}
%------------------------------------------------------------------------------%

%------------------------------------------------------------------------------%
\begin{example}
  Determine se a integral
  \(\int_{-1}^{0} \frac{1}{x^{\sfrac{2}{3}}}dx\)
  é convergente ou divergente.
  \solution
  Mais uma vez devemos nos atentar que a integral em questão é imprópria tipo 1,
  uma vez que o integrando possui uma descontinuidade em $x=0$.
  Uma vez que essa descontinuidade acontece no extremo direito do intervalo
  $[-1, 0]$, precisamos usar a parte 1 da Definição \ref{def:integral-impropria-2}.
  Assim,
  \begin{align*}
    \int_{-1}^{0} \frac{1}{x^{\sfrac{2}{3}}}dx
    & = \lim_{t \to 0^{-}}
        \int_{-1}^{t}
        \frac{1}{x^{\sfrac{2}{3}}}dx  \\
    & \stackTFCD
      \lim_{t \to 0^{-}}
      3x^{\sfrac{1}{3}} \evalat{-1}{t} \\
    & = \lim_{t \to 0^{-}}
      3 \left[ t^{\sfrac{1}{3}} - (-1)^{\sfrac{1}{3}} \right] \\
    & = 3
  \end{align*}
  Portanto, a integral imprópria converge.
\end{example}
%------------------------------------------------------------------------------%

%------------------------------------------------------------------------------%
\begin{example}\label{ex:imp3}
  Calcule \(\int_{0}^3 \frac{1}{x-1} dx\), se possível.
  \solution
  Primeiramente, é importante observar que a integral em questão é
  imprópria-tipo 2, uma vez que $\sfrac{1}{x-1}$ possui uma descontinuidade
  em $x=1$. Como essa descontinuidade ocorre dentro do intervalo $[0, 3]$,
  devemos utilizar a parte 3 da Definição~\ref{def:integral-impropria-2}
  com $c=1$.
  \[
      \int_{0}^3 \frac{1}{x-1} dx
    = \int_{0}^1 \frac{1}{x-1} dx
    + \int_{1}^3 \frac{1}{x-1} dx
  \]
  onde
  \begin{align*}
  \int_{0}^1 \frac{1}{x-1} dx
  & = \lim_{t \to 1^{-}}\int_{0}^t \frac{1}{x-1} dx  \\
  & \stackTFCD \lim_{t \to 1^{-}} \ln\abs{x-1} \evalat{0}{t}  \\
  & = \lim_{t \to 1^{-}}\big(\ln\abs{t-1} - \ln\abs{-1} \big) \\
  & = \lim_{t \to 1^{-}} \ln(1-t) =-\infty
  \end{align*}
  pois $1-t \to 0^+$, quando $t \to 1^{-}$.
  Portanto, como
  \[\int_{0}^1 \frac{1}{x-1} dx\] é divergente, então
  \[\int_{0}^3 \frac{1}{x-1} dx\] é divergente e não precisamos calcular
  \[\int_{1}^3 \frac{1}{x-1} dx \,.\]
\end{example}

Note que se não tivéssemos atentados que na função do exemplo~\ref{ex:imp3},
poderíamos ter feito \textbf{erroneamente} o seguinte cálculo
\[
  \int_{0}^3 \frac{1}{x-1} dx
  = \ln\abs{t-1} \evalat{0}{3}
  = \ln 2 - \ln 1
  = \ln 2 \,\,
\]
no entanto, esse procedimento está incorreto, uma vez que a integral
imprópria deve ser calculada em termos de limite. Isso nos leva a
uma importante lição: a partir de agora, sempre que for solicitado
calcular
\[
  \int_{a}^{b} f(x)dx \,,
\]
é fundamental decidir se a integral
é ordinária ou imprópria.

%------------------------------------------------------------------------------%
